% =====================================================================
% ejercicio05.tex
% Responsabilidad única: comparar divisiones por reducción de Gini,
% asignar clases en hojas, elegir una división de regresión por suma de
% errores cuadrados y discutir extrapolación y control de complejidad.
% Cifras verificadas con verificacion/teoria.py.
% =====================================================================
\subsection{Ejercicio 5: Árboles de decisión, división, predicción y complejidad}

\paragraph{La situación.} Un árbol de decisión predice haciendo preguntas
de sí o no sobre las variables, por ejemplo \emph{¿el equipo tiene más de
3 años?}, y cada respuesta manda a la observación por una rama distinta.
Cada pregunta parte un grupo de observaciones (un \emph{nodo}) en dos
grupos \emph{hijos}, y los grupos finales, donde ya no se pregunta más, se
llaman \emph{hojas}. En clasificación, la hoja predice la clase que más se
repite en ella; en regresión, el promedio de $y$ en ella. Para decidir qué
pregunta hacer, el árbol busca la que deje a los hijos lo más
\emph{puros} posible, es decir, con observaciones de una sola clase.

Para medir qué tan mezclado está un nodo usamos la \emph{impureza de
Gini}, $G=1-p_0^2-p_1^2$, donde $p_0$ y $p_1$ son las proporciones de
observaciones de clase $0$ y de clase $1$ en el nodo. Si el nodo es puro
(todo de una clase), $G=1-1^2-0^2=0$; si está mitad y mitad,
$G=1-0.25-0.25=0.5$, el máximo posible con dos clases. Una forma de leerla
es como la probabilidad de que, al sacar dos observaciones al azar del
nodo (devolviendo la primera), sean de clases distintas.

En este ejercicio un nodo tiene $12$ observaciones, $6$ de clase $0$ y $6$
de clase $1$, y comparamos dos preguntas candidatas. Escribimos cada hijo
con sus conteos $(n_0,n_1)$: la división $A$ produce hijos $(5,1)$ y
$(1,5)$, y la $B$ produce $(4,0)$ y $(2,6)$.

\begin{enumerate}

\item \textbf{Impureza Gini y reducción ponderada.}
\emph{Paso 1: impureza del padre.} Con $6$ de cada clase,
$p_0=p_1=\frac{6}{12}=\frac12$, y
\[
    G_{\text{padre}} = 1 - \bigl(\tfrac{1}{2}\bigr)^2 - \bigl(\tfrac{1}{2}\bigr)^2 = 1-\tfrac14-\tfrac14 = 0.5 .
\]

\emph{Paso 2: impureza de cada hijo.} Para cada hijo calculamos sus
proporciones y aplicamos la fórmula:
\begin{align*}
    (5,1):&\quad p_0=\tfrac56,\ p_1=\tfrac16, &
    G&=1-\tfrac{25}{36}-\tfrac{1}{36}=\tfrac{10}{36}=\tfrac{5}{18}\approx0.2778,\\
    (1,5):&\quad p_0=\tfrac16,\ p_1=\tfrac56, &
    G&=\tfrac{5}{18}\approx0.2778\quad\text{(simétrico al anterior)},\\
    (4,0):&\quad p_0=1,\ p_1=0, &
    G&=1-1-0=0\quad\text{(puro)},\\
    (2,6):&\quad p_0=\tfrac28=\tfrac14,\ p_1=\tfrac68=\tfrac34, &
    G&=1-\tfrac{1}{16}-\tfrac{9}{16}=\tfrac{6}{16}=\tfrac38=0.375 .
\end{align*}

\emph{Paso 3: reducción ponderada.} La mejora de una división es la
impureza del padre menos el promedio de las impurezas de los hijos,
ponderado por la fracción de observaciones que se lleva cada uno
($n_L$ y $n_R$ son los tamaños de los hijos izquierdo y derecho, y $n=12$):
\[
    \Delta G = G_{\text{padre}} - \frac{n_L}{n}G_L - \frac{n_R}{n}G_R .
\]
Para $A$, cada hijo tiene $6$ de las $12$ observaciones:
$\Delta G_A=0.5-\frac{6}{12}\cdot\frac{5}{18}-\frac{6}{12}\cdot\frac{5}{18}
=\frac12-\frac{5}{18}=\frac{9}{18}-\frac{5}{18}=\frac{2}{9}\approx0.2222$.
Para $B$, los hijos tienen $4$ y $8$ observaciones:
$\Delta G_B=0.5-\frac{4}{12}\cdot0-\frac{8}{12}\cdot\frac38
=\frac12-\frac14=0.25$. La Tabla~\ref{tab:ej5-gini} resume las cuentas.

\begin{table}[H]
\centering
\begin{tabular}{ccccccc}
\toprule
División & Hijo & $(n_0,n_1)$ & $n_h/n$ & $G_h$ & $\sum_h \frac{n_h}{n}G_h$ & $\Delta G$ \\
\midrule
\multirow{2}{*}{$A$}
  & izq. & $(5,1)$ & $6/12$ & $\frac{5}{18}\approx0.2778$
  & \multirow{2}{*}{$\frac{5}{18}\approx0.2778$} & \multirow{2}{*}{$\frac{2}{9}\approx0.2222$} \\
  & der. & $(1,5)$ & $6/12$ & $\frac{5}{18}\approx0.2778$ & & \\
\midrule
\multirow{2}{*}{$B$}
  & izq. & $(4,0)$ & $4/12$ & $0$
  & \multirow{2}{*}{$\frac14=0.2500$} & \multirow{2}{*}{$\frac14=\mathbf{0.2500}$} \\
  & der. & $(2,6)$ & $8/12$ & $\frac38=0.3750$ & & \\
\bottomrule
\end{tabular}
\caption{Impureza Gini de cada hijo y reducción ponderada por división;
por ejemplo, $G_{(5,1)}=1-\frac{25}{36}-\frac{1}{36}=\frac{5}{18}$ y
$G_{(2,6)}=1-\frac{4}{64}-\frac{36}{64}=\frac{3}{8}$, de modo que en $B$
la impureza ponderada es $\frac{4}{12}\cdot0+\frac{8}{12}\cdot\frac38=\frac14$.}
\label{tab:ej5-gini}
\end{table}

\emph{Paso 4: elección.} Elegimos la \textbf{división $B$}, con mayor
reducción de impureza ($0.25>0.2222$). A primera vista nos parecía mejor
$A$ por sus dos hijos simétricos, pero $B$ aísla un grupo completamente
puro de $4$ observaciones de clase $0$ (Figura~\ref{fig:ej5-arboles}), y
eso vale más que dejar dos grupos moderadamente mezclados.

\emph{Paso 5: por qué ponderar por el tamaño.} Lo que queremos reducir es
la mezcla que sufren las \emph{observaciones}, no la de los nodos como si
fueran unidades, así que un hijo con $8$ de las $12$ observaciones debe
pesar el doble que uno con $4$. Si promediáramos sin pesos, $B$ se
calificaría con $(0+0.375)/2=0.1875$, es decir, una reducción de $0.3125$ en
lugar de $0.25$: estaríamos premiando de más la pureza de un hijo pequeño.
Llevado al extremo, cualquier división que aislara una sola observación
(un hijo con $G=0$) parecería excelente sin haber mejorado casi nada la
clasificación del resto. Además, el promedio ponderado tiene una lectura
natural: es la impureza que en promedio encuentra una observación tomada
al azar del nodo después de la división.

\begin{figure}[H]
\centering
\newcommand{\ejcNodoGini}[6]{% #1 nombre, #2 x, #3 y, #4 n0, #5 n1, #6 texto de G
    \node[draw=gray!60, rounded corners=2pt, fill=white,
          minimum width=2.3cm, minimum height=1.45cm, inner sep=0pt]
          (#1) at (#2,#3) {};
    \node[font=\small, anchor=north] at ($(#1.north)+(0,-0.06)$) {$(#4,#5)$};
    \node[font=\footnotesize, anchor=north] at ($(#1.north)+(0,-0.48)$) {$G=#6$};
    \pgfmathsetmacro{\ejcFr}{1.9*#4/(#4+#5)}
    \fill[classblue] ($(#1.south west)+(0.2,0.15)$) rectangle ++(\ejcFr,0.22);
    \fill[classred] ($(#1.south west)+(0.2,0.15)+(\ejcFr,0)$)
        rectangle ($(#1.south west)+(2.1,0.37)$);
}
\begin{tikzpicture}[arista/.style={-{Stealth[length=2mm]}, gray!70!black}]
% --- División A ---
\node[font=\small\bfseries] at (0,1.25) {División $A$};
\ejcNodoGini{PA}{0}{0}{6}{6}{0.5}
\ejcNodoGini{LA}{-1.55}{-2.3}{5}{1}{5/18}
\ejcNodoGini{RA}{1.55}{-2.3}{1}{5}{5/18}
\draw[arista] (PA.south) -- (LA.north)
    node[midway, above left, inner sep=1.5pt, font=\footnotesize] {$6/12$};
\draw[arista] (PA.south) -- (RA.north)
    node[midway, above right, inner sep=1.5pt, font=\footnotesize] {$6/12$};
\node[font=\small] at (0,-3.65)
    {$\Delta G=\tfrac12-\tfrac{5}{18}=\tfrac29\approx0.2222$};
% --- División B ---
\node[font=\small\bfseries, text=classgreen!70!black] at (7.2,1.25)
    {División $B$ (elegida)};
\ejcNodoGini{PB}{7.2}{0}{6}{6}{0.5}
\ejcNodoGini{LB}{5.65}{-2.3}{4}{0}{0}
\ejcNodoGini{RB}{8.75}{-2.3}{2}{6}{3/8}
\draw[arista] (PB.south) -- (LB.north)
    node[midway, above left, inner sep=1.5pt, font=\footnotesize] {$4/12$};
\draw[arista] (PB.south) -- (RB.north)
    node[midway, above right, inner sep=1.5pt, font=\footnotesize] {$8/12$};
\node[font=\small, draw=classgreen, thick, rounded corners=2pt,
      inner sep=3pt] at (7.2,-3.65)
    {$\Delta G=\tfrac12-\tfrac14=\mathbf{0.25}$};
\draw[classgreen, thick, dashed, rounded corners=4pt]
    (4.15,1.65) rectangle (10.25,-4.15);
% --- Leyenda ---
\fill[classblue] (2.35,-4.75) rectangle ++(0.3,0.2);
\node[font=\footnotesize, anchor=west] at (2.7,-4.65) {clase $0$};
\fill[classred] (4.25,-4.75) rectangle ++(0.3,0.2);
\node[font=\footnotesize, anchor=west] at (4.6,-4.65) {clase $1$};
\end{tikzpicture}
\caption{Divisiones candidatas $A$ y $B$ del nodo $(6,6)$, donde cada
barra muestra la proporción de clase $0$ (azul) y de clase $1$ (rojo)
en el nodo y cada arista el peso $n_h/n$ del hijo; aunque $A$ produce
dos hijos simétricos y moderadamente impuros, $B$ aísla un hijo puro
que concentra un tercio de las observaciones, de modo que logra la
mayor reducción ponderada ($0.25$ frente a $0.2222$) y es la división
que elegimos.}
\label{fig:ej5-arboles}
\end{figure}

\item \textbf{Clases y probabilidades en las hojas.}
En la división $B$, cada hijo es una hoja. La probabilidad estimada de
clase $1$ en una hoja es simplemente la proporción de observaciones de
clase $1$ que cayeron en ella:
\begin{itemize}
    \item Hoja izquierda $(4,0)$: clase mayoritaria $0$ y
    $\hat P(Y=1)=0/4=0$.
    \item Hoja derecha $(2,6)$: clase mayoritaria $1$ y
    $\hat P(Y=1)=6/8=0.75$.
\end{itemize}
Usar la clase mayoritaria equivale a usar el umbral $t=0.5$. Con el umbral
$t=0.8$ la hoja izquierda sigue en clase $0$ ($0<0.8$) y la hoja derecha
\emph{también} pasa a clase $0$, pues $0.75<0.8$. Así \textbf{ambas hojas
predicen clase $0$} y el árbol deja de predecir la clase positiva en
absoluto. Esto muestra algo propio de los árboles: sus probabilidades solo
pueden tomar unos pocos valores (uno por hoja), así que mover el umbral
entre dos de esos valores no cambia ninguna predicción, y cuando cruza uno
de ellos cambia una hoja entera de golpe.

\item \textbf{División de regresión por SSE.}
Ahora la respuesta es numérica: $x=(1,\dots,6)$ y $y=(2,3,7,8,9,12)$. Cada
hoja predice el promedio de sus $y$, y medimos la calidad de una división
con la suma de errores al cuadrado (SSE) de las dos hojas.

\emph{División $x\le2.5$.} Quedan las hojas $\{2,3\}$ (los $x=1,2$) y
$\{7,8,9,12\}$ (los $x=3,\dots,6$):
\begin{align*}
    \text{izquierda: } & \bar y=\tfrac{2+3}{2}=2.5, &
    SSE&=(2-2.5)^2+(3-2.5)^2=0.25+0.25=0.5,\\
    \text{derecha: } & \bar y=\tfrac{7+8+9+12}{4}=9, &
    SSE&=(-2)^2+(-1)^2+0^2+3^2=4+1+0+9=14,
\end{align*}
para un total de $0.5+14=\mathbf{14.5}$.

\emph{División $x\le4.5$.} Quedan las hojas $\{2,3,7,8\}$ y $\{9,12\}$:
\begin{align*}
    \text{izquierda: } & \bar y=\tfrac{2+3+7+8}{4}=5, &
    SSE&=(-3)^2+(-2)^2+2^2+3^2=9+4+4+9=26,\\
    \text{derecha: } & \bar y=\tfrac{9+12}{2}=10.5, &
    SSE&=(-1.5)^2+1.5^2=2.25+2.25=4.5,
\end{align*}
para un total de $26+4.5=30.5$.

\emph{Elección y predicción.} Elegimos $x\le2.5$, que baja el SSE del nodo
sin dividir (que es $70.8333$, alrededor de su promedio $41/6\approx6.8333$)
a $14.5$, frente a $30.5$ de la otra. Por curiosidad revisamos también los
otros cortes posibles ($1.5$, $3.5$ y $5.5$ dan $42.8$, $22.6667$ y $38.8$),
y ninguno le gana a $2.5$. Como $x=7>2.5$, cae en la hoja derecha y
predecimos $\hat y(7)=\mathbf{9}$.

\begin{figure}[H]
\centering
\begin{tikzpicture}
\begin{axis}[
    width=0.66\textwidth, height=5.6cm,
    xmin=0, xmax=8.5, ymin=0, ymax=14,
    xlabel={$x$}, ylabel={$y$},
    legend pos=south east, legend cell align=left,
    legend style={font=\small}, tick label style={font=\small},
]
\addplot[only marks, mark=*, mark size=2.4pt, color=classblue]
    coordinates {(1,2) (2,3) (3,7) (4,8) (5,9) (6,12)};
\addlegendentry{datos}
\addplot[very thick, color=classred, const plot]
    coordinates {(0,2.5) (2.5,9) (8.5,9)};
\addlegendentry{árbol $x\le2.5$}
\addplot[domain=0:8.5, dashed, color=gray] {-0.0667 + 1.9714*x};
\addlegendentry{tendencia lineal (MCO)}
\addplot[only marks, mark=o, mark size=2.6pt, color=black] coordinates {(7,9)};
\draw[gray!70, densely dotted] (axis cs:6,0) -- (axis cs:6,14);
\end{axis}
\end{tikzpicture}
\caption{Predicción constante a trozos del árbol con la división
$x\le2.5$; a la derecha de $x=6$ (línea punteada, último valor
observado) el árbol mantiene $\hat y=9$ mientras la tendencia lineal
sigue creciendo. El círculo vacío marca $\hat y(7)=9$.}
\end{figure}

\item \textbf{La media minimiza el SSE; no extrapolación.}
\emph{Paso 1: plantear.} En una hoja $R$ con $|R|$ observaciones (las
barras indican cuántas hay), el árbol predice un mismo número $c$ para
todas. El error de esa hoja es $f(c)=\sum_{i\in R}(y_i-c)^2$, donde
$i\in R$ quiere decir \emph{las observaciones que están en la hoja}.
Queremos el $c$ que lo hace mínimo.

\emph{Paso 2: derivar e igualar a cero.} La derivada de $(y_i-c)^2$
respecto a $c$ es $2(y_i-c)\cdot(-1)$, así que
\[
    f'(c) = -2\sum_{i\in R}(y_i-c) = 0
    \;\Longrightarrow\; \sum_{i\in R}y_i - |R|\,c = 0
    \;\Longrightarrow\; c^* = \frac{1}{|R|}\sum_{i\in R}y_i = \bar y_R ,
\]
es decir, el promedio de la hoja.

\emph{Paso 3: es un mínimo.} La segunda derivada es $f''(c)=2|R|>0$, de
modo que $f$ es una parábola que abre hacia arriba y $\bar y_R$ es su único
mínimo. Otra forma de verlo, sin derivadas, es sumar y restar $\bar y_R$
dentro del cuadrado:
\[
    f(c)=\sum_{i\in R}\bigl[(y_i-\bar y_R)+(\bar y_R-c)\bigr]^2
    =\sum_{i\in R}(y_i-\bar y_R)^2
     +2(\bar y_R-c)\underbrace{\sum_{i\in R}(y_i-\bar y_R)}_{=0}
     +|R|(\bar y_R-c)^2 .
\]
El primer término no depende de $c$ y el último es un cuadrado que solo
vale cero cuando $c=\bar y_R$; cualquier otro valor de $c$ suma error.

\emph{Paso 4: por qué no extrapola.} Como cada hoja predice una
constante, la hoja que recibe todos los $x$ mayores que el último corte
cubre el intervalo $(2.5,\infty)$ completo, y entonces
$\hat y(7)=\hat y(20)=\hat y(1000)=9$. El árbol \textbf{no puede
extrapolar} una tendencia: sus predicciones siempre son promedios de datos
de entrenamiento, así que nunca salen del rango de esos promedios (aquí
$[2.5,9]$). Y eso que los datos sugieren un crecimiento sostenido: la
recta de mínimos cuadrados sobre estos seis puntos sube aproximadamente
$1.97$ por cada unidad de $x$, y seguiría subiendo más allá de $x=6$.

\item \textbf{\texttt{max\_depth}, \texttt{min\_samples\_leaf} y sobreajuste.}
Si dejamos que un árbol siga preguntando sin límite, termina con una hoja
para cada observación y se aprende de memoria los datos de entrenamiento,
incluido su ruido. Los dos hiperparámetros le ponen freno de formas
distintas:
\begin{itemize}
    \item \texttt{max\_depth} (profundidad máxima) limita cuántas preguntas
    seguidas puede haber desde la raíz hasta una hoja. Como cada nivel a lo
    sumo duplica el número de hojas, un árbol de profundidad $d$ tiene a lo
    sumo $2^{d}$ hojas, y eso limita qué tan fina puede ser la partición e
    impide llegar a aislar observaciones individuales.
    \item \texttt{min\_samples\_leaf} (mínimo de observaciones por hoja)
    exige que cada hoja conserve al menos ese número de observaciones de
    entrenamiento. Así cada predicción (promedio o proporción) se calcula
    con más datos y es menos sensible al ruido de unas pocas observaciones.
\end{itemize}
Si aumentamos demasiado la profundidad, esperamos este patrón: el error de
\textbf{entrenamiento} baja siempre, hasta llegar prácticamente a cero,
porque un árbol sin límites puede dar a cada observación su propia hoja;
el error de \textbf{validación}, en cambio, primero baja (el árbol está
capturando estructura real) y luego se estanca o sube (empieza a memorizar
ruido). Se forma la típica curva en U para validación y una brecha cada vez
mayor entre ambos errores, que es la señal de sobreajuste que buscamos
controlar en los ejercicios 7 y 9 al seleccionar estos hiperparámetros con
el conjunto de validación.

\end{enumerate}
